\documentclass[a4paper,11pt]{ctexbook}
\usepackage{xcolor}
\usepackage{graphicx}
\usepackage[top=2cm, bottom=2cm, left=2.2cm, right=2.2cm]{geometry}
\usepackage{array}
\usepackage{hyperref}
\usepackage{booktabs}
\usepackage{amsmath}
\usepackage{amssymb}
\usepackage{amsthm}
\usepackage{framed}
\usepackage{marginnote}
\usepackage{tcolorbox}

\definecolor{themeblue}{RGB}{0,80,150}
\definecolor{scenario}{RGB}{180,50,30}
\definecolor{quant}{RGB}{30,120,70}
\definecolor{sidebarbg}{RGB}{245,248,252}
\definecolor{summarybg}{RGB}{230,245,230}
\definecolor{discussbg}{RGB}{255,245,230}

\newenvironment{scenario}{
  \vspace{0.4cm}
  \noindent\begin{tcolorbox}[colback=scenario!5!white,colframe=scenario!70!black,title=\textcolor{scenario}{\textbf{情景引入}}]
  \vspace{0.1cm}
}{\end{tcolorbox}\vspace{0.3cm}}

\newenvironment{qualitative}{
  \vspace{0.3cm}
  \noindent\begin{tcolorbox}[colback=themeblue!5!white,colframe=themeblue!60!black,title=\textcolor{themeblue}{\textbf{定性分析}}]
  \vspace{0.1cm}
}{\end{tcolorbox}\vspace{0.3cm}}

\newenvironment{quantitative}{
  \vspace{0.3cm}
  \noindent\begin{tcolorbox}[colback=quant!5!white,colframe=quant!60!black,title=\textcolor{quant}{\textbf{定量分析}}]
  \vspace{0.1cm}
}{\end{tcolorbox}\vspace{0.3cm}}

\newenvironment{lessonsummary}{
  \vspace{0.3cm}
  \noindent\begin{tcolorbox}[colback=summarybg,colframe=green!50!black,title=\textbf{\textcolor{green!50!black}{总结}}]
  \vspace{0.1cm}
}{\end{tcolorbox}\vspace{0.3cm}}

\newenvironment{discuss}{
  \vspace{0.3cm}
  \noindent\begin{tcolorbox}[colback=discussbg,colframe=orange!60!black,title=\textbf{\textcolor{orange!60!black}{讨论题}}]
  \vspace{0.1cm}
}{\end{tcolorbox}\vspace{0.3cm}}

\newenvironment{exercises}{
  \vspace{0.3cm}
  \noindent\begin{tcolorbox}[colback=blue!3!white,colframe=blue!50!black,title=\textbf{\textcolor{blue!50!black}{课后练习}}]
  \vspace{0.1cm}
}{\end{tcolorbox}\vspace{0.3cm}}

\newcommand{\xuehai}[1]{\marginnote{\footnotesize\textbf{\textcolor{themeblue}{学海拾遗}}\\ #1}}
\newcommand{\xuehaiinline}[1]{
  \vspace{0.2cm}
  \noindent
  \fcolorbox{themeblue!40}{sidebarbg}{
    \begin{minipage}{0.92\textwidth}
    \textbf{\textcolor{themeblue}{【学海拾遗】}}\\ #1
    \end{minipage}}
  \vspace{0.2cm}}

\newenvironment{figplaceholder}[1]{
  \vspace{0.2cm}
  \begin{center}
  \fbox{\begin{minipage}{0.82\textwidth}\centering\textcolor{blue!70!black}{\textbf{【插图位置】}} \\[0.2cm] #1 \end{minipage}}
  \end{center}
  \vspace{0.2cm}}

\linespread{1.4}

\title{\textcolor{themeblue}{\Huge\textbf{机械工程基础}}}
\subtitle{\textcolor{gray}{E01·第一卷\quad 共六课}}
\author{}
\date{}

\begin{document}
\maketitle
\thispagestyle{empty}
\tableofcontents
\newpage

% ============================================================
\chapter{静力学——力为什么能平衡？}
% ============================================================

\begin{scenario}
\textbf{塔吊倒塌的瞬间}

2021年北京某工地，一台塔吊在吊装预制构件时突然倒塌——起重臂折断，配重块砸向地面。调查发现：塔吊基础的地脚螺栓有2颗松动，导致应力重新分布后超过极限。一个螺栓的松动，摧毁了整台塔吊。静力学要回答的根本问题就是——在什么条件下，一个结构或机械能保持不动？"不动"意味着所有力相互抵消——这就是平衡。
\end{scenario}

\begin{qualitative}
\textbf{平衡方程——力系简化的核心}

如果一个物体处于静止状态或匀速直线运动状态，那么作用在它上面的所有力必须满足两个条件：

\begin{enumerate}
  \item 所有力的矢量和为零——物体不会加速
  \item 所有力对任意一点的力矩之和为零——物体不会转动
\end{enumerate}

对于平面力系，这就是三个标量方程：
\[
\sum F_x = 0, \quad \sum F_y = 0, \quad \sum M_O = 0
\]

"静定" vs "超静定"：如果未知力的数量恰好等于独立平衡方程的数量——问题是静定的（可解）。如果未知力更多——是超静定的，需要补充变形协调方程。

\textbf{自由体受力图（Free Body Diagram, FBD）：}把研究对象从周围环境中"隔离"出来，画出所有作用在该物体上的外力。这是静力学分析的第一步，也是最关键的一步——任何一个力的遗漏或方向错误，都会导致整个分析错误。
\end{qualitative}

\begin{quantitative}
\textbf{桁架分析——两种方法}

\textbf{节点法：}逐个取桁架的节点为研究对象，对每个节点列平衡方程。适合求所有杆件的内力。

\textbf{截面法：}假想一个截面将桁架切断，取其中一部分列平衡方程。适合求某几根特定杆件的内力。

\textbf{案例：}一个简单的简支桁架，跨度24米，上弦受均布荷载。给定每根杆件的截面面积$A$和材料的许用应力$[\sigma]$，求最大安全荷载$P_{max}$。

\textbf{摩擦力：}静摩擦 $f_s \le \mu_s N$，动摩擦 $f_k = \mu_k N$。摩擦角 $\phi = \arctan\mu$。在机械设计中，螺栓连接的预紧力计算和自锁条件分析都依赖摩擦角。螺旋千斤顶的自锁条件：螺纹升角$\lambda < \arctan f$。
\end{quantitative}

\begin{lessonsummary}
\begin{enumerate}
  \item 平衡方程是所有机械和结构静力分析的起点——$\sum F=0$、$\sum M=0$
  \item 自由体受力图（FBD）是静力学分析的第一步——画错了FBD一切白做
  \item 桁架分析的节点法和截面法分别适用于整体分析和特定杆件分析
  \item 静定系统和超静定系统的区分是分析的第一步——后者需要补充变形协调条件
\end{enumerate}
\end{lessonsummary}

\begin{discuss}
\begin{enumerate}
  \item 为什么塔吊的底座有那么多地脚螺栓而不是只有几颗粗大的？螺栓越多，应力越分散——但同时也意味着任何一个螺栓失效都会引起应力重新分布。这是一个"冗余度"和"单点故障"之间的工程权衡。
  \item "超静定"结构在实际工程中比比皆是——比如桥梁的连续梁支撑、飞机的机翼连接。为什么工程师要设计超静定结构而不是简单的静定结构？——超静定结构在其中一根支撑失效后不会直接倒塌，有"安全冗余"。
  \item 摩擦在机械设计中是"好的"还是"坏的"？——螺栓连接（需要摩擦来锁紧）是好的；轴承（需要减少摩擦来提高效率）是坏的。工程师的工作就是在这个"好"和"坏"之间找到平衡。
\end{enumerate}
\end{discuss}

\begin{exercises}
\begin{enumerate}
  \item 一个简支梁跨度为6米，在4米处作用一个8kN的集中力。画出梁的FBD，计算两个支座的约束反力。
  \item 一个平面桁架有7个节点、11根杆件。验证是否满足静定条件 $m + 3 = 2j$（$m$=杆件数，$j$=节点数）。如果满足，用节点法计算所有杆件的内力。
  \item 一块重500N的铁块放在水平地面上，静摩擦系数0.3。需要多大的水平推力才能让它开始滑动？如果推力方向与水平面成30°角呢？
\end{enumerate}
\end{exercises}

\xuehaiinline{
\textbf{工程师的自由体图——为什么老师总是强调"画FBD"？}FBD的本质是把复杂的真实世界简化为一个清晰的分析模型。真实世界中，一根梁上可能有几十个力——自重、外加荷载、支撑反力、风载、温度应力……FBD要求你"选择"哪些力是相关的、应该被画出来，哪些力是次要的、可以忽略。这个"选择"的决策能力——区分"重要"和"不重要"——是一个工程师从新手到专家的分水岭。\\
\textbf{超静定系统的"力法"：}求解超静定结构时，关键步骤是释放多余的约束→得到一个静定基→然后补充"变形协调条件"（在原来的约束位置，位移必须为零）。这个思想——把一个"太紧"的问题先"松绑"、再"找补回来"——是工程分析的核心思维方式，不限于结构力学。
}

\newpage
% ============================================================
\chapter{材料力学——一个零件能承受多大的力？}
% ============================================================

\begin{scenario}
\textbf{泰坦尼克号的钢材——为什么冰冷的海水让它变脆了？}

1912年4月15日，泰坦尼克号撞上冰山后沉没。1985年残骸被发现后，科学家对打捞上来的钢板进行了测试——发现泰坦尼克的钢材在冰冷的海水中（约-2°C）发生了\textbf{脆性转变}。同一块钢板在室温下有良好的延展性（可以弯曲变形），但在0°C以下变得像玻璃一样脆——受到冲击时会直接断裂而不是变形。这种"温度诱导的脆化"在材料力学中被称为\textbf{韧脆转变温度}——它是材料科学和力学交叉的核心问题。
\end{scenario}

\begin{qualitative}
\textbf{应力与应变——材料对力的"回答"}

当你对一根金属棒施加拉力时，金属棒会变长——拉力除以截面面积是\textbf{应力} $\sigma$，伸长量除以原长是\textbf{应变} $\epsilon$。

在弹性范围内，应力和应变满足胡克定律：
\[
\sigma = E\epsilon
\]

$E$是\textbf{弹性模量（杨氏模量）}——它是材料"刚度"的度量。钢的$E\approx 200$GPa，铝合金约70GPa，聚合物约1-5GPa。$E$越高的材料，在同样应力下变形越小。

\textbf{应力-应变曲线的四个阶段：}
\begin{enumerate}
  \item \textbf{弹性区：}可恢复的变形，遵循胡克定律
  \item \textbf{屈服：}应力不再增加但应变继续增加——材料"屈服"了
  \item \textbf{强化：}屈服后材料重新获得承载能力——应力继续上升
  \item \textbf{颈缩：}局部截面缩小→最终断裂
\end{enumerate}

\textbf{安全系数：}$n = \sigma_{limit} / \sigma_{working}$。飞机零件安全系数通常1.5，桥梁2.0，电梯钢丝绳10+。
\end{qualitative}

\begin{quantitative}
\textbf{拉压、扭转与弯曲}

\textbf{轴向拉压：}$\sigma = F/A$，$\Delta L = FL/(EA)$。

\textbf{扭转：}$\tau = Tr/J$（圆轴），$\theta = TL/(GJ)$。$G$为剪切模量，$J$为极惯性矩。

\textbf{弯曲：}梁的弯曲正应力 $\sigma = My/I$，$I$为截面惯性矩。对于矩形截面$I = bh^3/12$，对于圆截面$I = \pi d^4/64$。

\textbf{应力集中：}当截面突然变化（孔、槽、台阶）时，局部应力远大于名义应力：
\[
\sigma_{max} = K_t \times \sigma_{nom}
\]
$K_t$为应力集中系数——对于圆孔，$K_t \approx 3$。

\textbf{案例：}一根直径20mm的钢杆受拉力50kN。计算应力：$\sigma = 50000/(\pi \times 0.01^2) \approx 159$MPa。$E=200$GPa，应变量$\epsilon = 159/200000 = 0.0008$，伸长量$\Delta L = 0.08\%$。如果杆上有一个直径5mm的横向通孔，局部应力峰值 $= 159 \times 3 \approx 477$MPa——已经超过了大多数结构钢的屈服强度。
\end{quantitative}

\begin{lessonsummary}
\begin{enumerate}
  \item 应力是单位面积上的力（强度问题），应变是相对变形（刚度问题）
  \item 胡克定律 $\sigma = E\epsilon$ 是线弹性力学的基石——但只在弹性范围内成立
  \item 应力集中是导致大部分机械疲劳失效的根源——设计时应该避免尖锐的截面突变
  \item 安全系数是"不知道的不知道"的工程保险——所有设计中都必须考虑
\end{enumerate}
\end{lessonsummary}

\begin{discuss}
\begin{enumerate}
  \item "安全系数越大越好"吗？——过大的安全系数意味着浪费材料、增加重量和成本。在飞机设计中（轻量化是核心目标），安全系数通常只有1.5。而在电梯中（人命关天），钢丝绳的安全系数高达10以上。一个系统应该在哪里"省"、在哪里"余"？这是工程中最重要的权衡决策之一。
  \item 2021年虎门大桥发生涡振（在风中像"水蛇"一样扭动）——大桥的应力水平实际上远低于设计极限（安全系数充足），但振动幅度足以引起公众恐慌。这说明工程中"安全"和"可接受"是两个不同的标准。
  \item 为什么用高强度钢制造的零件不一定"更好"？——高强度钢的屈服强度高，但断裂韧性可能更低（更脆）。在需要承受冲击或低温的场景中，中等强度但高韧性的材料可能是更好的选择。
\end{enumerate}
\end{discuss}

\begin{exercises}
\begin{enumerate}
  \item 一根直径10mm的铝合金杆（$E=70$GPa，$\sigma_y=250$MPa），长度2米，承受10kN拉力。计算应力、应变和伸长量。该杆件是否进入塑性？
  \item 一个矩形截面梁（$b=50$mm，$h=100$mm）受最大弯矩10kN·m。计算梁上下表面的最大弯曲应力。
  \item 一根直径30mm的圆轴承受500N·m的扭矩。计算最大剪切应力。如果轴的材料许用剪应力80MPa，安全系数是多少？
\end{enumerate}
\end{exercises}

\xuehaiinline{
\textbf{泰坦尼克号钢材的确切分析（MIT 2003）：}1996年和1998年的两次打捞行动获取了泰坦尼克号的钢板样本。MIT和加拿大纽芬兰纪念大学的测试发现：泰坦尼克的钢材含硫量较高（与现代钢相比约3倍），这使得钢材的韧脆转变温度在约+30°C——意味着在冰海温度（-2°C）下，钢材处于完全的脆性状态。撞击冰山时，钢板不是像现代造船钢那样"变形吸收能量"——而是"直接裂开"。\\
\textbf{不仅如此——铁达尼号的铆钉也存在问题：}研究发现部分铆钉含大量矿渣夹杂物（冶金质量差），在低温下同样变脆。撞击冰山时，铆钉头断裂→钢板接缝张开→海水涌入。一个宏观的灾难，根源在于微观的材料缺陷。\\
\textbf{今天的造船钢材：}必须满足-40°C（北极航线）或-60°C（极地科考船）的低温韧性要求。标准测试方法——夏比冲击试验（Charpy V-notch test）——测量标准试件在特定温度下被摆锤打断时吸收的能量。
}

\newpage
% ============================================================
\chapter{动力学——力和运动为什么互为因果？}
% ============================================================

\begin{scenario}
\textbf{内燃机——如何用"爆炸"产生旋转运动？}

汽车发动机的汽缸里，汽油和空气的混合气体被火花塞点燃——瞬间燃烧膨胀，推动活塞向下运动。但车轮需要的是"旋转"而不是"上下"。发动机把活塞的直线运动通过曲柄连杆机构变成曲轴的旋转运动——这就是动力机械最经典的转换：往复直线→旋转。每秒重复几十次的微小爆炸，驱动着一辆两吨重的汽车以120km/h在高速公路上巡航。
\end{scenario}

\begin{qualitative}
\textbf{牛顿定律在机械系统中的应用}

动力学研究的是"力如何改变运动状态"。它包含两个互补的问题：
\begin{enumerate}
  \item 已知力→求加速度（正问题：$F=ma$）
  \item 已知运动→求力（逆问题：从加速度反推需要的力）
\end{enumerate}

对于旋转运动，牛顿第二定律变为：
\[
M = I\alpha
\]

其中$M$是力矩（扭矩），$I$是转动惯量，$\alpha$是角加速度。

\textbf{四杆机构与曲柄滑块：}最常见的机械运动转换机构。曲柄做整周旋转，通过连杆带动滑块做往复直线运动。发动机的活塞-连杆-曲轴就是一个标准的曲柄滑块机构。
\end{qualitative}

\begin{quantitative}
\textbf{转动惯量的计算}

对于简单几何体的转动惯量：
\begin{itemize}
  \item 细圆环绕轴线：$I = mr^2$
  \item 圆盘绕轴线：$I = \frac12 mr^2$
  \item 细杆绕一端：$I = \frac13 mL^2$
  \item 细杆绕中点：$I = \frac{1}{12} mL^2$
\end{itemize}

\textbf{平行轴定理：}$I = I_{cm} + md^2$——如果旋转轴不穿过质心，需要加$md^2$项。

\textbf{曲柄滑块机构的运动学：}曲柄角速度$\omega$，连杆长度$L$，曲柄半径$R$，滑块位移：
\[
x = R\cos\omega t + L\sqrt{1 - (R/L)^2\sin^2\omega t}
\]
当$R/L < 0.3$时，可近似为$x \approx R\cos\omega t + const$。
\end{quantitative}

\begin{lessonsummary}
\begin{enumerate}
  \item $F=ma$ 和 $M=I\alpha$ 是机械动力学的两个基本方程——覆盖了平动和转动
  \item 曲柄-连杆（曲柄滑块机构）是把旋转运动变成直线运动（或反之）的最常见机构
  \item 转动惯量的计算是旋转运动分析的基础——形状和旋转轴位置决定了$I$
  \item 机械系统中90\%的运动是平动和转动的组合——分析的关键是选择正确的"自由度"
\end{enumerate}
\end{lessonsummary}

\begin{discuss}
\begin{enumerate}
  \item 为什么内燃机的活塞需要"活塞环"？——活塞环的作用是密封（防止燃气泄漏到曲轴箱）和传热。但活塞环也增加了摩擦损耗——约占总摩擦损失的40\%。取消活塞环→降低摩擦→提高效率，但密封问题怎么解决？——这就是"无活塞环发动机"的研究挑战。
  \item "往复惯性力"——活塞上下运动产生的惯性力是发动机振动的主要来源。直列四缸发动机的二级往复惯性力无法完全平衡，所以四缸机天然地比六缸机振动更大。这就是为什么豪华车用六缸、八缸——不完全是动力需求，而是平顺性需求。
  \item 电动汽车的电机不需要曲柄连杆——转子直接产生旋转运动。这是否意味着电动车的机械系统比燃油车更简单？——是的。电动车传动系统的零件数量约为燃油车的1/5。但简单不一定意味着更可靠——永磁电机的稀土磁铁在高温下会退磁、功率电子器件（IGBT/SiC模块）的可靠性仍然是挑战。
\end{enumerate}
\end{discuss}

\begin{exercises}
\begin{enumerate}
  \item 一个飞轮的转动惯量$I=10$kg·m$^2$，在5秒内从静止加速到角速度100rad/s，需要施加多大的恒定扭矩？
  \item 一个曲柄滑块机构：曲柄半径$R=50$mm，连杆长度$L=200$mm，曲柄转速3000rpm。求滑块的最大速度和最大加速度。
  \item 一个四缸发动机的曲柄夹角为180°（每缸间的点火间隔相等）。画出各缸活塞的往复惯性力，分析二阶惯性力的情况。
\end{enumerate}
\end{exercises}

\xuehaiinline{
\textbf{为什么"平衡"对发动机如此重要？}一个单缸发动机的活塞在上下运动时，会产生方向交替变化的惯性力——这就是振动源。在1900年代，汽车工程师们发现即使发动机被牢固地固定在车架上，整个车身仍然剧烈抖动。\\
\textbf{解决思路：}（1）多缸布置——各缸的惯性力相互抵消一部分；（2）加平衡轴——专门制造一个反向旋转的配重轴来抵消二阶惯性力；（3）采用完全平衡的设计——水平对置发动机（Boxer engine）的活塞相向而行，惯性力在水平方向完全抵消——这就是保时捷和斯巴鲁使用水平对置发动机的原因。\\
\textbf{一个历史趣事：}1920年代的工程师在计算发动机振动时，忽略了二阶惯性力（频率是曲轴转速的两倍），认为"反正是高频，无所谓"。结果装上车的发动机在某个特定转速下剧烈共振——因为二阶激励频率恰好和车架的固有频率重合了。从此汽车工程师把"固有频率分析"作为设计的必选项。
}

\newpage
% ============================================================
\chapter{能量与功率——机器效率的极限在哪？}
% ============================================================

\begin{scenario}
\textbf{永动机的诱惑与不可能}

从13世纪开始，无数人试图制造"永动机"——不需要外部能量输入、可以永远对外做功的机器。所有尝试都失败了。为什么？因为能量守恒定律说——能量不能凭空产生，也不能凭空消失。热力学第二定律进一步说——不可能从单一热源吸取热量并将其完全转化为功而不产生其他影响。一个机器可以转换能量（从燃料到运动），但转换效率永远不可能达到100\%。机械工程的一个重要使命就是：让这个转化效率尽可能接近热力学的极限。
\end{scenario}

\begin{qualitative}
\textbf{功、能量和效率}

\textbf{功：}力在位移上的累积——$W = \vec{F} \cdot \vec{d}$。

\textbf{动能：}平动 $E_k = \frac12 mv^2$，转动 $E_k = \frac12 I\omega^2$。

\textbf{势能：}重力势能 $E_p = mgh$，弹性势能 $E_p = \frac12 kx^2$。

\textbf{功率：}$P = Fv$（平动），$P = M\omega$（转动）。

\textbf{机械效率：}
\[
\eta = \frac{P_{out}}{P_{in}} = \frac{E_{out}}{E_{in}}
\]

齿轮传动的效率通常在90-99\%之间；液压系统60-85\%；内燃机25-45\%；电动机85-95\%。
\end{qualitative}

\begin{quantitative}
\textbf{卡诺效率——热机的极限}

任何热机（内燃机、蒸汽轮机、燃气轮机）都工作在高温热源（$T_H$）和低温热源（$T_L$）之间。卡诺热机的效率——所有热机效率的理论上限：
\[
\eta_{Carnot} = 1 - \frac{T_L}{T_H}
\]

所有温度应当使用绝对温度（开尔文）。\\
\textbf{案例：}汽油发动机的燃烧温度约2500K、排气温度约800K。卡诺极限效率 $\eta = 1 - 800/2500 = 68\%$。但实际汽油机的最高效率只有约35-40\%——原因：不完全燃烧、摩擦损失、泵气损失、冷却损失等。

\textbf{能量守恒在机械系统中的应用：}一个电机—减速器—负载系统。电机输出功率$P_m$，减速器效率$\eta_g$，负载获得的功率$P_l = P_m \times \eta_g$。损失的功率转化为热能——这就是为什么变速箱需要冷却。

\textbf{飞轮储能：}$E = \frac12 I\omega^2$——飞轮可以储存旋转动能，在需要时释放。F1赛车的KERS（动能回收系统）就使用飞轮储能。
\end{quantitative}

\begin{lessonsummary}
\begin{enumerate}
  \item 能量既不能创生也不能消灭——一切效率低于100\%的机器都会产生热量
  \item 卡诺效率是所有热机的理论极限——它只取决于高温端和低温端的温度
  \item 机械系统的能量损失来源：摩擦（转化为热）、噪音（振动能）、热传导
  \item 提高效率的工程方向：减少摩擦（轴承/润滑）、提高工作温度（发动机）、回收废能（KERS/涡轮增压）
\end{enumerate}
\end{lessonsummary}

\begin{discuss}
\begin{enumerate}
  \item 电动车的能量效率（约85-90\%从电池到车轮）远高于燃油车（约20\%从汽油到车轮）。但为什么人类社会花了一百年时间才"想到"用电力驱动取代内燃机？——因为电池的能量密度（约250Wh/kg）远远低于汽油（约12000Wh/kg）。1kg汽油包含的能量需要48kg电池才能储存。所以能源效率只是技术选择的一个维度——能量密度、功率密度、成本、基础设施可兼容性同样重要。
  \item "废热"到底"废"了多少？——一座1000MW的燃煤发电厂，按照卡诺极限（$\eta \approx 40\%$），约有1500MW的热量（以蒸汽/热水/烟气的形式）被"浪费"了。如果能把这部分废热用于区域供暖——就是热电联产（CHP），综合能量利用率可以从40\%提升到80\%以上。
  \item 制氢——用可再生电力电解水制氢（Power-to-Gas），再通过氢燃料电池发电。这个链路的每一步都有能量损失：电解（70-80\%效率）→压缩运输（80-90\%）→燃料电池（50-60\%）→电机（90\%）——总的"电→电"效率大约只有25-35\%。直接用电（电→电）效率约90\%。这算是一种"浪费"吗？——从效率的角度：是的。但氢气可以被储存（解决可再生能源的间歇性问题），而电力不能被大规模储存——所以效率不是唯一的标准。
\end{enumerate}
\end{discuss}

\begin{exercises}
\begin{enumerate}
  \item 一台汽油发动机输出功率100kW，效率30\%。它每秒钟消耗多少焦耳的燃料能量？如果改用效率40\%的柴油机，同样输出功率，燃料能量消耗减少多少？
  \item 一个电机（功率50kW，效率90\%）通过减速器（效率95\%）驱动一个负载。负载获得的功率是多少？减速器中每秒钟损失多少能量？
  \item 一个飞轮的转动惯量50kg·m$^2$，转速3000rpm时可以储存多少能量？如果转速降到1500rpm，剩余能量是多少？
\end{enumerate}
\end{exercises}

\xuehaiinline{
\textbf{永动机的著名骗局——Redheffer的"永动机"（1812）：}Charles Redheffer在美国费城展示了一台"永动机"——一个大型轮子在没有任何可见动力源的情况下持续转动。工程师们百思不得其解，直到有人发现机器背后的墙上有异常——打开后发现墙后有一个老人在用手摇曲柄驱动着这台"永动机"。Redheffer逃跑了——这个故事说明：当某样东西违反能量守恒时，要么你理解错了，要么有人在骗你。\\
\textbf{今天还有一种合法的"永动机"——在地球以外的太空中：}如果空间中的卫星不受到大气阻力（在足够高的轨道），它可以靠角动量守恒"永久"旋转——这不是永动机，因为没有对外做功。一旦航天器进行姿态调整（喷气/飞轮），角动量守恒被破坏——需要消耗燃料或电能。\\
\textbf{热泵——看起来"效率超过100\%"的机器：}热泵不产生热量——它从低温热源（室外/地下水/土壤）吸热、转移到高温热源（室内）。热泵的COP（性能系数，即输出热量/输入电能）可以达到3-4——不是违反热力学定律，而是因为它把"搬运热量"也算了进去。
}

\newpage
% ============================================================
\chapter{机械设计——公差、配合与表面粗糙度}
% ============================================================

\begin{scenario}
\textbf{一个"对不上"的螺栓孔}

小明在工厂实习，他的任务是加工一个法兰盘上的四个螺栓孔。他按图样精心加工——完成后拿到另一个车间去配对时发现：螺栓孔和对应的法兰上的孔怎么也"对不上"——差了0.5mm。他"严格按照图样"加工的，为什么会错？问题出在"基准"——他的基准和对方车间的基准不一致。但这个"0.5mm"在工程上真的很要紧吗？——如果螺栓直径比孔径小3mm，0.5mm的偏移是没问题的；但如果螺栓直径和孔径的间隙只有0.2mm——这0.5mm的偏移就导装不上。
\end{scenario}

\begin{qualitative}
\textbf{公差——零件的"身份证"}

在机械设计中，不可能制造出"完全精确"的零件——每个零件都有误差。公差就是"允许的误差范围"。它决定了：
\begin{itemize}
  \item 两个零件是否能够装配在一起（装配性）
  \item 装配后的精度和性能（功能）
  \item 制造成本（公差越紧，成本呈指数级上升）
\end{itemize}

\textbf{配合的三种类型：}
\begin{itemize}
  \item \textbf{间隙配合：}孔的最小直径大于轴的最大直径——可以自由滑动
  \item \textbf{过盈配合：}孔的最大直径小于轴的最小直径——需要压入或热装
  \item \textbf{过渡配合：}可能有间隙也可能有过盈——取决于实际尺寸落在公差范围的位置
\end{itemize}

\textbf{基孔制与基轴制：}优先采用基孔制——孔更难加工，所以以孔为基准（H7/h6等）。
\end{qualitative}

\begin{quantitative}
\textbf{公差等级与粗糙度}

\textbf{IT公差等级：}IT01到IT18共20级——IT6是高精度（磨削）、IT9中等精度（精车）、IT12粗精度（粗车）。IT6的尺寸公差大约是$10i$，其中$i$是公差单位（$i \approx 0.45\sqrt[3]{D} + 0.001D$）。

\textbf{表面粗糙度：}$Ra$（轮廓算术平均偏差）是最常用的参数。Ra=0.1μm（镜面磨削）、Ra=1.6μm（精车）、Ra=12.5μm（粗铣）。

\textbf{尺寸链：}多个零件的公差累积到最终装配间隙。$T_{tot} = \sum T_i$（极值法）或 $T_{tot} = \sqrt{\sum T_i^2}$（平方和根法，假设独立正态分布）。

\textbf{案例：}三个零件叠起来总长度100mm，每个零件的公差$\pm 0.1$mm。极值法：总公差$\pm 0.3$mm。平方和法：总公差$\pm 0.17$mm。设计时的选择决定了装配的成功率。
\end{quantitative}

\begin{lessonsummary}
\begin{enumerate}
  \item 公差是设计的"精确程度"——太松（功能不合格）、太紧（成本指数级上升）
  \item 配合类型（间隙/过盈/过渡）决定了装配方式和功能
  \item 尺寸链分析是多零件装配中保证"装得上"的关键工具
  \item 表面粗糙度影响摩擦、磨损、密封和疲劳寿命——每提高一级Ra精度，制造成本约翻倍
\end{enumerate}
\end{lessonsummary}

\begin{discuss}
\begin{enumerate}
  \item 为什么"6σ"质量管理在制造业如此重要？——零件在公差范围内时被认为是"合格"的，但如果一个尺寸靠近公差上限而另一个靠近公差下限——装配时可能"装不上"。6σ的"过程能力指数$C_{pk}$"要求实际分布范围只有公差带的一半——即使用"正态分布"的极端尾部也不会超出公差范围。
  \item "三坐标测量机（CMM）"可以测量零件的实际尺寸到微米级精度。但测量结果和"图纸设计值"之间的差距——是零件不合格、还是测量本身有误差？测量误差（测量不确定度）本身也必须被量化——通常要求测量不确定度不大于公差的10\%。
  \item 中国制造业的"精度瓶颈"——为什么高端机床（五轴加工中心、高精度磨床）仍然大量依赖进口？精密机床的关键在于"导轨的直线度"和"主轴的旋转精度"——二者都依赖极高精度的机械加工。一台中端数控车床可能需要加工80多个零件，每个的公差在IT6-IT7级别——这些公差的实现必须依赖更精密的"母机"（机床的机床）。中国的高精度母机仍然落后于日本和德国。
\end{enumerate}
\end{discuss}

\begin{exercises}
\begin{enumerate}
  \item 一个轴和孔的配合要求是$\phi 50H7/g6$。查公差表（或计算）：轴和孔的极限尺寸分别是多少？配合类型是什么？
  \item 三个零件叠装设计总长100$\pm 0.2$mm，每个零件的加工公差按平方和法分配，每个零件的公差是多少？
  \item 一个零件需要$Ra 0.8\mu m$的表面粗糙度，目前只能达到$Ra 3.2\mu m$。从粗车到磨削，成本翻倍大约需要多少倍？如果该零件的年产量10万个，每个零件多花5元钱加工费——总成本增加50万元。这个"更光滑的表面"真的值50万吗？
\end{enumerate}
\end{exercises}

\xuehaiinline{
\textbf{世界上最精密的机械加工——VLSI光刻机的镜头：}荷兰ASML公司的EUV（极紫外）光刻机，其反射镜系统是"人类制造的最光滑的表面"——表面粗糙度$Ra$达到0.02nm（纳微米），相当于把地球表面磨平到小于1mm的起伏。镜子的形状精度要求是几十个皮米。这些镜子在制造过程中耗时数月、用离子束逐点修正。没有这些"极端精度"的镜子——芯片就无法刻出7nm、5nm的线宽——现代电子产品就不存在了。\\
\textbf{零件的互换性——一个200年前的想法：}在18世纪之前，每一支枪都是"一对一"制造的——一个枪管的零件只能装在该枪管上，坏了就必须找制造者重新配。1798年，Eli Whitney（美国轧棉机的发明者）向美国政府展示了一个概念——用同样的模具和夹具制造1000支步枪的零件，然后随机挑出零件装配——全部完美匹配。这是"互换性制造"的诞生——没有它，就没有流水线、没有量产、没有现代工业。
}

\newpage
% ============================================================
\chapter{北京在地——中国制造业与机械工程}
% ============================================================

\begin{scenario}
\textbf{亦庄——北京"制造"了什么？}

很多人以为北京只有科技和金融——但北京经济技术开发区（亦庄）聚集了超过200家世界500强企业的制造基地。京东方（BOE）的8.5代液晶面板生产线——全厂自动化物流、机械臂搬运玻璃基板（2.5米×2.2米、厚度仅0.5mm）不能有任何颤抖。北京奔驰的工厂——每90秒有一辆整车下线，焊接车间500多个机器人自动完成6000多个焊点。亦庄代表了"中国制造2025"在北京的实践——不是"廉价制造"，而是"精密制造"和"智能制造"。
\end{scenario}

\begin{qualitative}
\textbf{中国机械工程行业的现状}

中国是全球最大的制造业国家——制造业增加值占全球约30\%（超过美国+日本+德国之和）。但大并不等于强：

\textbf{优势：}
\begin{itemize}
  \item 产业链完整——从螺丝钉到高铁的全链条制造能力
  \item 规模效应——同样的产品，中国的制造成本通常比发达国家低20-40\%
  \item 基础设施——高速铁路/港口/电力的制造和建设能力全球领先
\end{itemize}

\textbf{短板：}
\begin{itemize}
  \item 高端装备自主化率低——高端数控机床约70\%依赖进口，工业传感器约80\%依赖进口
  \item 工业软件——CAD/CAE/CAM/PLM基本被西门子/达索/PTC垄断
  \item 基础材料——高端轴承钢、模具钢、碳纤维等仍然依赖进口
\end{itemize}
\end{qualitative}

\begin{quantitative}
\textbf{北京机械工程相关的产业数据}

\begin{itemize}
  \item 亦庄经济技术开发区GDP占北京市约25\%——是北京最大的单一经济功能区
  \item 北京奔驰2023年产量约60万辆——年产值约3000亿元
  \item 京东方2023年营收约1700亿元——全球LCD面板出货量第一
  \item 中关村科技园有超过40家与智能制造相关的上市公司
  \item 北京航空航天大学、清华大学、北京理工大学的机械工程学科均为A+评级
\end{itemize}
\end{quantitative}

\begin{lessonsummary}
\begin{enumerate}
  \item 北京不是只有服务业——亦庄是我国北方最大的先进制造业基地之一
  \item 中国机械工程行业的"大而不强"体现在高端装备和基础工业软件两大领域
  \item 智能制造（工业4.0）的核心——不是"机器换人"，而是"数据驱动决策"
  \item 从机械工程的角度，"卡脖子"技术（光刻机、高端传感器、工业软件）本质上是"精度和可靠性"的问题
\end{enumerate}
\end{lessonsummary}

\begin{discuss}
\begin{enumerate}
  \item "卡脖子"的实质——以高端轴承钢为例，瑞典SKF、德国FAG、日本NSK垄断了全球高端轴承。中国能生产轴承（量大、便宜），但高铁轴承、精密机床主轴轴承的寿命和可靠性差30-50\%。这不是"做不出来"而是"做不精"——材料纯度+热处理工艺+表面改性+装配精度的综合差距。
  \item 国产替代的路径：Copy → Analyze → Improve → Create。中国制造业在许多领域走过了"复制"和"分析"阶段，正在进入"改进"和"创造"阶段——但最难的"从99\%到99.99\%"的过程，需要的基础研究投入和时间，可能是前几个阶段的总和。
  \item 为什么高端工业软件这么难替代？——不是因为算法（很多算法都是公开发表的），而是因为（1）软件产业生态（插件/模板/用户社区/培训体系）不是一两年能建起来的；（2）汽车和航空行业几十年的经验和验证已经固化到了这些软件中——从CAD模型到制造数据到质检报告的整个流程都是基于这些软件建立的。"替代"意味着推倒重来——这个成本没有人愿意承担。
\end{enumerate}
\end{discuss}

\begin{exercises}
\begin{enumerate}
  \item 选择一家北京制造业企业（如京东方/北京奔驰/北汽集团），通过网络或公开信息了解它的主要产品和制造工艺。写一份300字的"这家公司用了哪些机械工程技术"的分析。
  \item 网上搜索"高端数控机床国产化率"的最新数据（2024-2025年），分析近5年国产化率的变化趋势。哪些领域进步最快？哪些最慢？为什么？
  \item 如果让你选一个"卡脖子"技术方向作为未来大学专业的方向——高端轴承、精密传感器、工业软件、精密测量仪器——你会选哪个？为什么？在学术检索平台上查一下这个方向最近5年的国家自然科学基金资助情况，看看"该方向的投入"和"被卡的程度"是否匹配。
\end{enumerate}
\end{exercises}

\xuehaiinline{
\textbf{北京制造业的几个"隐形冠军"：}
\begin{itemize}
  \item \textbf{北方华创（NAURA）：}中国最大的半导体设备制造商之一，总部在亦庄——刻蚀机、薄膜沉积设备在全球市场份额约3-5\%（正在快速追赶应用材料/东京电子/Lam Research）
  \item \textbf{北京精雕：}生产高精度数控雕刻机和加工中心——主轴精度达到0.1微米级，在精密模具加工领域国内领先
  \item \textbf{中科三环：}全球最大的钕铁硼永磁材料供应商之一——中国稀土永磁材料产量占全球约90\%，这本身就是一种"挟稀土以令诸侯"的工业实力
\end{itemize}
\textbf{看制造业的一个"隐藏指标"——机床的数控化率：}中国每年消费的机床中，数控化率约30\%（日本的约90\%）。这意味着中国制造业中仍然有大量老旧的手动和非数控机床——不是说"造不出来"，而是中小制造企业没有资金和动机去更新设备。这个"存量"的改造任务，本身就是未来10年最大的机械工程市场。
}

\end{document}
